% Pipeline: entanglement mass into inherited Newton. Not derived Poisson.
\documentclass[11pt]{article}
\usepackage[margin=1.05in]{geometry}
\usepackage{amsmath,amssymb,amsthm}
\usepackage[colorlinks=true,linkcolor=blue,citecolor=blue,urlcolor=blue]{hyperref}

\newtheorem*{admission}{Admission}

\title{\textbf{Entanglement Reads the Mass.\\
Inherited Newton Does $1/r^2$.}}
\author{Robert W.\ McGwier\\
\small CTO, Cohere Technology Group\\
\small GitHub: \texttt{n4hy}}
\date{15 August 2026 --- Version 1}

\begin{document}
\maketitle

\begin{abstract}
\noindent
Paper~47 used $\kappa$ to weigh and locate. This note
feeds $M_{\mathrm{hat}}$ into the locked M9.2 DST Poisson
solver.

$N=12$, $R=3$, packet at $(6,6,6)$. $\kappa=0.968$ from
the even well-inside split. On the odd half,
$M_{\mathrm{hat}}=\mathrm{median}(\delta S/\kappa)$ matches
enclosed energy to $1.7\times 10^{-8}$. The high-$\delta S$
intersection is the true site.

That mass, placed as a compact blob at the centre of a
Dirichlet cube ($n=65$, $L=1$, $G=1$), is attractive
$GM/r^2$ to $3.0\%$, $2.1\%$, $1.9\%$ at
$0.30L$, $0.35L$, $0.40L$. Slope $-2.037$. C1 PASS.

Auditor, $N=10$, own $M_{\mathrm{hat}}$: mass CONFIRMED
($0.18\%$). C1 CONFIRMED. The C1 residuals are the same
number on both masses because Poisson is linear.

Entanglement supplied the mass. Einstein still did
$1/r^2$. Not a derivation of Poisson. Not $8\pi G$ from
$\kappa$. Not $1/4G$. Not FGHMV. Not de~Sitter.
\end{abstract}

\begin{admission}
I ran the pipeline. The force gate passed because I put a
positive mass into Poisson. I did not write that
entanglement derived $1/r^2$.
\end{admission}

\end{document}
